\section{Cấu hình và khả năng tái lập}

Phụ lục này tóm tắt các cấu hình có ảnh hưởng trực tiếp đến kết quả. Báo cáo
giữ nguyên dữ liệu đầu vào và không thay đổi số liệu trong quá trình biên tập.
Quy trình tái lập gồm kiểm tra schema và ranh giới train/test, tạo rolling
features, huấn luyện theo engine, đánh giá trên test và đối chiếu các bảng kết
quả sau khi ghi.

\subsection{Môi trường và cấu hình chính}

\begin{table}[H]
\centering
\caption{Các cấu hình chính của thực nghiệm.}
\label{tab:environment}
\begin{tabular}{ll@{\qquad}ll}
\toprule
Thành phần & Phiên bản/Giá trị & Thành phần & Phiên bản/Giá trị \\\midrule
Python & 3.12.x & Pandas & 2.3.3 \\
NumPy & 2.3.3 & SciPy & 1.16.2 \\
scikit-learn & 1.9.1 & statsmodels & 0.14.5 \\
PyArrow & 23.0.1 & PySpark & 4.0.1 \\
Window & 30 cycles & GroupKFold & 5 folds \\
Seed & 42 & Bootstrap metric & 10.000 lần \\
\bottomrule
\end{tabular}
\end{table}

Bootstrap cuối lấy mẫu lại engine có hoàn lại và giữ số lần xuất hiện của mỗi
engine khi tính mean engine-level MAE. Wilcoxon dùng 100.000 permutation
resamples với Benjamini--Hochberg trên 21 sensors. Nhánh PCA--ARIMA fit
transformer trên 90 engine train và chỉ transform 10 engine được đánh giá; mỗi
chuỗi dùng 70\% đầu để fit và 30\% cuối để forecast. Các giá trị này mô tả đúng
cấu hình đã chạy, không phải mặc định tối ưu cho mọi subset C-MAPSS.

\subsection{Nguyên tắc tái lập}

Các phép kiểm tra được thực hiện theo thứ tự: xác nhận schema và cycle
continuity; tách dữ liệu train/test; tạo nhãn RUL và rolling features; chạy
EDA, paired analysis và PCA--ARIMA trên train; thực hiện grouped
cross-validation; ghi final predictions rồi mới ghép test labels. Kết quả được
đối chiếu lại theo schema, khóa, số dòng và giá trị số sau mỗi bước ghi dữ liệu.

Quy trình này kiểm tra khả năng tái lập của lần chạy chính trên FD001 với cấu
hình đã nêu. Ma trận bổ sung trong \filepath{outputs/cross\_dataset/matrix.json}
kiểm tra riêng transfer từ FD001 sang FD002--FD004 với window 10 cycles. Kết quả
không chứng minh scalability trên cluster, khả năng khái quát tới động cơ thực
tế hoặc tính phù hợp của metric cho quyết định bảo trì.

\subsection{Nguồn đối chiếu kết quả}

\begin{table}[H]
\centering
\caption{Đối chiếu các nhóm kết quả trong báo cáo.}
\label{tab:result-map}
\small
\begin{tabular}{L{4.2cm}L{5.2cm}L{5.1cm}}
\toprule
Nhóm kết quả & Nội dung đối chiếu & Kiểm tra liên quan \\\midrule
Chất lượng dữ liệu & Schema, khóa và độ liên tục của cycle & Row count, kiểu số hữu hạn và khóa duy nhất \\
Rolling features & Đặc trưng từ 30 cycles gần nhất & Đủ window, cùng engine và không dùng tương lai \\
Model RUL & Cross-validation và metric test & GroupKFold, sample weighting và test-label join sau prediction \\
Suy luận sensor & Chênh lệch đầu/cuối vòng đời & Cặp một--một, bootstrap, Wilcoxon và FDR \\
Health Index & PCA--ARIMA và baseline cùng horizon & Temporal holdout, convergence, residual diagnostics và coverage \\
\bottomrule
\end{tabular}
\end{table}
